\section{Introduction}

A common problem in large content-serving platforms is the need to assign accurate searchable tags to content. These tags are important because they help users discover relevant content through recommendation systems, filtering, and manual search. In many cases, tags may be created manually by content creators or automatically by the platform, but both approaches can introduce errors. As a result, building systems that can automatically predict useful and accurate tags from textual information is an important natural language processing task. \\

In this project, we focus on the task of classifying movies into genres using their textual descriptions. This problem is relevant to streaming platforms and movie databases, where genre labels help organize content and improve user experience. Our approach takes a movie's plot description as input and predicts a single primary genre as output. Although movies may naturally belong to multiple genres, we simplify the task to single-label classification in order to make the prediction setting more consistent and easier to evaluate.\\

This task is challenging because movie plot descriptions are often short, vague, and sometimes misleading without additional context. Descriptions may emphasize only one part of the story while leaving out details that are important for identifying genre. In addition, some genres have overlapping characteristics, which makes classification more difficult even for human readers. Despite these challenges, textual descriptions still contain useful patterns that can be learned by machine learning and neural network models.\\

The goal of this project is to design and evaluate multiple approaches for predicting a movie's genre from its plot description. We compare simple baselines, feature-based models such as logistic regression, neural models including recurrent neural networks (RNNs), and pretrained transformer-based approaches such as BERT and zero-shot classification. This allows us to analyze how different modeling strategies perform on the same task and understand the trade-offs between simpler and more complex methods.